% Generated draft: requires mathematical and source review.
\documentclass[11pt]{article}
\usepackage[margin=25mm]{geometry}
\usepackage{fontspec}
\setmainfont{Latin Modern Roman}
\usepackage{amsmath,amssymb,mathtools,unicode-math}
\setmathfont{Latin Modern Math}
\usepackage{enumitem,needspace,xurl,hyperref}
\setlength{\parindent}{0pt}
\setlength{\parskip}{5pt}
\setlength{\emergencystretch}{3em}
\newcommand{\N}{\mathbb{N}}
\newcommand{\Z}{\mathbb{Z}}
\newcommand{\Q}{\mathbb{Q}}
\newcommand{\R}{\mathbb{R}}
\newcommand{\C}{\mathbb{C}}
\newcommand{\F}{\mathbb{F}}
\newcommand{\A}{\mathbb{A}}
\newcommand{\PP}{\mathbb{P}}
\newcommand{\E}{\mathbb{E}}
\newcommand{\eps}{\varepsilon}
\newcommand{\dd}{\,\mathrm{d}}
\begin{document}
\section*{UnsolvedMath ID: 30003412}
\textbf{Source problem code:} catalogue entry
\subsection{Definitions and mathematical statement}
For an uncountable cardinal \(\kappa\), a \(\kappa\)-tree is a partially ordered set \((T,<_{T})\) such that the strict predecessors of each node are well ordered by \(<_{T}\), the tree has height \(\kappa\), and every level has cardinality less than \(\kappa\). More explicitly, if \(\operatorname{ht}_{T}(t)\) is the order type of the predecessors of \(t\), then
\[
T_\xi=\{t\in T:\operatorname{ht}_{T}(t)=\xi\},\qquad
0<|T_\xi|<\kappa\quad(\xi<\kappa),
\]
and there are no nodes of height at least \(\kappa\). A cofinal branch is a linearly ordered subset whose node heights are unbounded in \(\kappa\). A \(\kappa\)-Aronszajn tree is a \(\kappa\)-tree with no cofinal branch. The tree property at \(\kappa\), denoted here by \(\operatorname{TP}(\kappa)\), asserts that every \(\kappa\)-tree has a cofinal branch.

Let \(\omega_1\) be the first uncountable ordinal and let \(\aleph_\beta\) denote the \(\beta\)-th infinite cardinal, all computed within the same model. Taking a limit ordinal to be nonzero and to have no immediate predecessor, the supplied problem asks whether the following theory is consistent:
\[
\mathsf{ZFC}+\left[\forall\alpha<\omega_1\,
\bigl(\alpha\text{ is a nonzero limit ordinal}\ \Longrightarrow\
\operatorname{TP}(\aleph_{\alpha+1})\bigr)\right].
\]
Thus the tree properties are required simultaneously in one model for every countable nonzero limit ordinal. Interpreting consistency over \(\mathsf{ZFC}\) is the usual ambient convention; the supplied record does not explicitly specify a base theory or a large-cardinal hypothesis for a relative consistency result. The nonzero convention for limit ordinals is also not explicit in the source record; including \(0\) would add the separate requirement \(\operatorname{TP}(\aleph_1)\).
\subsection{Short English statement}
Can there consistently be a single model of set theory in which, for every countable nonzero limit ordinal \(\alpha\), every tree of height \(\aleph_{\alpha+1}\) with levels of size less than \(\aleph_{\alpha+1}\) has a cofinal branch? The base theory and the convention excluding zero from limit ordinals require confirmation.
\subsection{Sources}
\begin{itemize}
\item Local UnsolvedMath registry, supplied complete source record: permanent ID \(30003412\), title ``Tree Property at Successors of Countable Limit Alephs,'' source problem code \texttt{catalogue entry}, collection ``Oberwolfach Reports Open Problems.'' The record gives proposed year 2017 and names no proposer.
\item Supplied report citation: \emph{Set Theory (2018)}, Oberwolfach Reports, pp.~527--589, DOI \texttt{10.4171/owr/2017/11}. \url{https://doi.org/10.4171/owr/2017/11}. The supplied title year is 2018, whereas the DOI contains 2017; this discrepancy has not been resolved.
\item Supplied literature source: ``OWR Set Theory report.'' \url{https://ems.press/content/serial-article-files/46674}.
\item Supplied literature source: ``Diagonal supercompact Radin forcing.'' \url{https://doi.org/10.1016/j.apal.2020.102828}.
\item Literature assessment dated exactly \texttt{2026-08-21}: the supplied curation labels the problem open and states that no forcing construction currently establishes the simultaneous absence of Aronszajn trees in the stated pattern, while known tree-property results cover restricted patterns. This is supplied dated curation, not an independently verified literature assessment. The background explicitly requests independent verification of time-sensitive status and cited claims.
\item The supplied statement status is \texttt{corrected\_verified}. Its verification note reports that the official Oberwolfach report was checked directly and that the clean statement preserves the target while restoring definitions and removing source-dependent wording. These are supplied provenance claims; no independent source inspection is asserted here.
\end{itemize}
\end{document}
